% !TeX root = main.tex
\chapter{Implementation Results and Discussion}\label{chap:4}
\textbf{Completion outline---not final results.} No Chapter Four draft was
included in the supplied report-review Markdown set. The structure below follows
the materials review; populate it from the actual, dated evidence records.

\section{Implementation Checkpoint and Evaluation Scope}
Identify the software revision, terminal build, datasets, participants and test
conditions used for each evaluation. Keep implemented, proposed and untested
functions distinct.
\section{Student Records and Document Workflows}
Present expected-versus-observed backend, browser, record-integrity and access
results. Report document upload, retrieval and revision outcomes only where
implemented and tested. Add the actual filing-task dataset and any comparable
manual baseline, with task counts and consistent timing definitions.
\section{Terminal Implementation and Assembly}
Add the actual circuit schematic, component values, verified connections,
supply measurements and assembly photographs. Put diagrams in folder/ and
photographs in plates/; label and cite each figure or plate in the text.
\section{Fingerprint Recognition Results}
Use the documented pilot's trial records. Separate whole-gallery identification
and selected-slot verification denominators, wrong identities, rejections,
processing errors and timing observations. Do not merge the later saved-probe
feasibility experiment into the pilot dataset.
\section{Reference Preservation and Pending Capture}
Present byte-integrity, restoration, capture persistence, USB transfer and pending
ingestion evidence separately. Pending ingestion alone is not completed
identity-to-attendance reconciliation.
\section{Integration Results and Limitations}
Report completed identity binding, roster, permission, timing, retry and duplicate
handling cases, identifying any remaining links explicitly. Compare findings to
the literature only where the evaluated conditions and metrics are comparable.
